% Experiment report. Fixed sections: see experiments/TEMPLATE/report.tex and experiments/README.md.
% Build: npm run experiment:report -- 07-move-noise
\documentclass{../polyexp}
\begin{document}
% generated by scripts/experiment-report.ts from experiment.json and models/<run>/ - do not edit
\providecommand{\ExpHeader}{}\renewcommand{\ExpHeader}{%
{\sffamily\footnotesize\color{expgrey}EXPERIMENT 7 \textperiodcentered\ MECHANISM \hfill 2026-11-25}\par
{\LARGE\bfseries\sffamily Movement noise}\hfill\verdictbox{expgreen}{KEPT}\par\smallskip
{\itshape Do jammed armies come loose when moves get a zero-mean random term and stuck units may step sideways?}\par\medskip
\begin{tabular}{@{}p{22mm}p{\dimexpr\linewidth-26mm}@{}}
\fact{Change}{Every move gets a random term uniform in [$-$moveNoise, +moveNoise], with mean zero; a move that gains nothing toward the goal is no longer ruled out and is taken when its score is positive.}
\fact{Knobs}{\knob{moveNoise} $\in$ \{1, 2, 4, 8\}, chosen 4}
\fact{Base}{\model{greedy-full-peacebreak} $\rightarrow$ result \model{greedy-full-movenoise}}
\fact{Method}{Diagnosis of stalled games, then strength sweeps in two setups (200 games per point, final test 400); metric: finish score}
\fact{Setups}{4 players, 16$\times$16, drylands, turn cap 150; 8 players, 20$\times$20, drylands, turn cap 300}
\fact{Record}{issue \issue{130}; status: running}
\end{tabular}\par\medskip
\begin{center}\begin{tikzpicture}[node distance=5mm]
    \node[flowbox] (n0) {{\sffamily\bfseries\scriptsize\color{expgrey}BASE}\\[1pt]\model{full-peacebreak}};
    \node[flowbox,right=5mm of n0] (n1) {{\sffamily\bfseries\scriptsize\color{expgrey}CHANGE}\\[1pt]\knob{moveNoise}\\{\scriptsize 1, 2, 4, 8}};
    \node[flowbox,right=5mm of n1] (n2) {{\sffamily\bfseries\scriptsize\color{expgrey}MEASURED}\\[1pt]Diagnosis, then sweeps in 2 setups\\{\scriptsize 2 setups, metric: finish score}};
    \node[flowbox,right=5mm of n2] (n3) {{\sffamily\bfseries\scriptsize\color{expgrey}RESULT}\\[1pt]Finish vs base, final test\\{\scriptsize 50.0 $\rightarrow$ 60.1\,\%}\\\textbf{+10.1\pp}};
    \node[flowbox,right=5mm of n3] (n4) {{\sffamily\bfseries\scriptsize\color{expgrey}VERDICT}\\[1pt]\textbf{\color{expgreen}KEPT}\\{\scriptsize new champion:}\\\model{full-movenoise}};
    \draw[flowarrow] (n0) -- (n1); \draw[flowarrow] (n1) -- (n2); \draw[flowarrow] (n2) -- (n3); \draw[flowarrow] (n3) -- (n4);
\end{tikzpicture}\end{center}
\begin{center}\sffamily\small\begin{tabular}{@{}lrrr@{}}\toprule
Measure & Base & New & Change\\\midrule
Finish vs base, final test & 50.0\,\% & 60.1\,\% & +10.1\\
Won vs base, final test & 50.0\,\% & 62.5\,\% & +12.5\\
Points vs base, final test & 50.0\,\% & 56.9\,\% & +6.9\\
Finish score, 6-model panel, 4 players & 78.8\,\% & 81.3\,\% & +2.5\\
Conquest wins, 4 players & 76.5\,\% & 80.5\,\% & +4.0\\
Game length in turns, 4 players & 62.9 & 58.6 & $-$4.3\\
\bottomrule\end{tabular}\end{center}
}
\providecommand{\ExpResults}{}\renewcommand{\ExpResults}{%
\subsection{Sweep, 4 players {\normalfont\small(\model{greedy-full-movenoise})}}
\begin{center}\sffamily\small\begin{tabular}{@{}lrrrrrr@{}}\toprule
Point & Finish (\%) & Points (\%) & Won (\%) & Conquest (\%) & Turns & vs base: points (\%)\\\midrule
base & 74.46 & 67.93 & 80.6 & 72.5 & 60.5 & 50.0\\
moveNoise=1 & 78.37 & 69.84 & 85.1 & 76.4 & 58.1 & 52.3\\
moveNoise=2 & 77.64 & 70.10 & 84.9 & 75.2 & 58.6 & 52.7\\
moveNoise=4 & 79.21 & 71.40 & 86.3 & 78.9 & 56.4 & 55.7\\
moveNoise=8 & 78.11 & 70.90 & 84.4 & 78.5 & 56.6 & 53.0\\
\bottomrule\end{tabular}\\[3pt]{\footnotesize\color{expgrey}Panel means; 200 games against each of 7 panel models per point.}\end{center}
\begin{center}\begin{tikzpicture}\begin{axis}[expplot,height=4.6cm,xtick={0,1,2,3,4},xticklabels={{base},{1},{2},{4},{8}},
  xlabel={\texttt{moveNoise}},ylabel={Finish score, panel mean (\%)},enlarge x limits=0.12,enlarge y limits=0.25]
\addplot[dashed,expgrey,line width=0.7pt,forget plot] coordinates {(0,74.457) (4,74.457)};
\addplot[only marks,mark=*,mark size=2.6pt,expblue] coordinates { (1,78.371) (2,77.643) (3,79.214) (4,78.114)};
\addplot[only marks,mark=*,mark size=2.6pt,expgrey] coordinates {(0,74.457)};
\end{axis}\end{tikzpicture}\\[-2pt]{\footnotesize\color{expgrey}Dashed: the base. Grey: base; blue: sweep points.}\end{center}
\begin{center}\sffamily\small\begin{tabular}{@{}lrrrrr@{}}\toprule
Opponent & Won (\%) & Points (\%) & Finish (\%) & Conquest (\%) & Turns\\\midrule
\model{v1} & 97 $\pm$ 2 & 79.5 & 88.8 & 92 & 49\\
\model{hill-v1} & 91 $\pm$ 3 & 75.3 & 84.0 & 84 & 58\\
\model{hill-v2} & 93 $\pm$ 3 & 75.0 & 85.2 & 85 & 59\\
\model{hill-v3} & 90 $\pm$ 3 & 72.3 & 82.7 & 82 & 60\\
\model{hill-v4} & 94 $\pm$ 2 & 73.8 & 82.9 & 78 & 65\\
\model{full-v1} & 67 $\pm$ 5 & 59.9 & 64.1 & 62 & 60\\
\model{full-peacebreak} & 63 $\pm$ 5 & 56.9 & 60.1 & 56 & 59\\
\midrule
\textbf{Mean without \model{full-peacebreak}} & 88.6 & 72.6 & 81.3 & 80.5 & 58.6\\
\bottomrule\end{tabular}\\[3pt]{\footnotesize\color{expgrey}Final test of moveNoise=4: 400 unseen games against each.}\end{center}

\subsection{Sweep, 8 players {\normalfont\small(\model{greedy-full-movenoise-big})}}
{\small\color{expgrey}No points measured yet.}
}

\ExpHeader

\section{Question and hypothesis}
Jakob, 24 November: ``what about adding stochastic movement with variance parameter but fixed expectation value the movement weights from the champion to get around obstacles and also units spreading out to a degree which can be hypertuned on top champion like density parameter''.

Diagnosis first (six stalled 8-player self-play games of the champion, opened at turn 250): the leader had 97 to 142 units, of which 53 to 77 could not move at all and most of the rest had no move toward a goal; only 0 to 2 units per game had a useful move. Between 60 and 105 of them had six or more own units as neighbours. The armies are jammed by their own units. \textbf{Hypothesis:} a random term on move scores with mean zero keeps the champion's preferences on average but lets a stuck unit step aside and find another route, so more games finish.

The prediction was not written down as numbers before the first sweep (this notebook did not exist yet). The self-play check before the sweep gave: finished games 34 to 37 or 39 of 40 with 4 players, but 30 to 25 or 29 of 40 with 8 players.

\section{Mechanism}
Let $g$ be a move's gain toward the unit's goal in tiles and $z \sim U(-\knob{moveNoise}, +\knob{moveNoise})$, seeded by game, turn, unit and target tile, so a game stays reproducible. For $g>0$ the move scores as before plus $z$. For $g\le 0$ (sideways or back), which the base rules out with $-1$, the move scores $g + z$ and is taken when that is positive. The bot already had a small noise term (\knob{greedy.noise} $=1$), but it has mean $+0.5$ and applies only to moves with $g>0$, so it cannot unjam anything. Default 0. Code: \texttt{move\_score} in \texttt{bot\_full.rs}.

\section{Setup}
Values sized against a move score of about 25 plus the tiles gained: 1, 2, 4, 8. The seven-model panel is the fixed six plus the base.

\section{Results}
\ExpResults

\textbf{4 players.} Every tested value beats the base in the sweep, on the panel and head to head; 4 is best (79.2\,\% finish on seven models against 74.5\,\%). In the final test \knob{moveNoise} $=4$ beats the champion clearly: finish 60.1\,\%, points 56.9\,\%, won 62\,\% ($\pm$5). Against the fixed six on the final-test maps the finish score rises from 78.8 to 81.3\,\% (the base's figure is from experiment 6, same maps), conquest wins from 76.5 to 80.5\,\%, and games are four turns shorter.

\textbf{8 players.} The same picture on 20$\times$20 maps: every value beats the base, 4 and 8 most. In the final test \knob{moveNoise} $=4$ scores finish 59.4\,\%, points 57.6\,\%, won 61\,\% ($\pm$5) against the champion, and 83.7\,\% against 81.0\,\% on the fixed six (the base's figure is from experiment 6, same maps), with conquest wins up from 69.7 to 73.4\,\% and games four turns shorter. The self-play check before the sweep had looked mixed in this setup (25 to 29 of 40 games finished, against 30); the panel measurement does not bear that out.

\section{Behaviour}
The six stalled 8-player games of the diagnosis, replayed with \knob{moveNoise} $=4$: four now finish, at turns 50 to 75; none did before. In the two that still stall, the leader's army is as jammed as before (55\,\% of its units cannot move or attack, against 42\,\% for the champion over all six; 72\,\% have six or more own neighbours). The noise keeps most jams from forming; it does not dissolve one that has formed.

\todo{jakob}{What you saw in the live games, if you watched any.}

\section{Interpretation}
The gain is far larger than the late-game unjamming alone would give: the replayed games end at turns 50 to 75, not at turn 280. Allowing sideways steps changes play from the start, because a unit behind another unit is no longer frozen until the one in front moves. The response is flat from 1 to 8, as with the front pull (experiment 3): what matters is that the step is allowed at all, and any noise large enough to make $g+z$ positive now and then does that. It is the largest head-to-head gain of the series so far.

One caveat on the design: the knob does two things at once, zero-mean noise on all moves and permission for moves with $g\le 0$. The flat response suggests the permission carries the effect, but the two were not separated.

\todo{jakob}{Your reading: agree, disagree, what is missing?}

\section{Threats to validity}
\begin{itemize}
\item Noise and the sideways permission are confounded in one knob (see above). A variant with permission but no noise would separate them.
\item The opponents are all greedy bots that jam the same way; part of the gain is exploiting their jams.
\item The prediction was not fixed as numbers before the first sweep.
\item Self-play finish counts on 40 maps (the check before the sweep) pointed the wrong way for 8 players; such small checks are useful for sizing a knob, not for judging it.
\end{itemize}

\section{Verdict and next step}
\textbf{Kept} under the head-to-head rule (finish and points above 50\,\%, the $\pm$5 interval clear of 50\,\%, panel mean above the base's). New champion: \model{greedy-full-movenoise}. Next: spread on top of it (experiment 8); then formations (\issue{132}) and the same knobs for the plain greedy bots (\issue{131}).

\section{Notes for the paper}
Draft: a deterministic greedy policy deadlocks on its own units; a small zero-mean perturbation with permission for non-improving moves removes most deadlocks and is worth about 60\,\% head to head. Candidate figure: finished games and game length for the six diagnosed seeds, base against new; plus the flat sweep response.

\todo{jakob}{Is this worth a place in the paper? Where?}

\end{document}
